% =============================================================================
% Section 6: Robustness Programme
% =============================================================================
\section{Robustness Programme}\label{sec:robustness}

This section subjects the ISI composite ranking to four distinct
perturbation exercises.  Each test isolates a different source of
methodological sensitivity: (i)~axis inclusion, (ii)~standardisation
choice, (iii)~weight specification, and (iv)~outlier treatment.  The
baseline is the ISI~v1.0 ranking (\cref{tab:ranking}).

\subsection{Leave-one-axis-out (LOO)}\label{sec:loo}

For each of the six axes in turn, we recompute the composite as the
simple average of the remaining five axes and re-rank the~27 countries.
The Spearman rank correlation~$\rho$ between the reduced ranking and
the baseline ranking measures the influence of the omitted axis.

% Table T9: LOO sensitivity summary
\begin{table}[htbp]
\centering
\caption{Leave-one-axis-out (LOO) rank sensitivity.  Spearman $\rho$ and Kendall
$\tau$ measure rank-order agreement between the baseline and each LOO ranking.
Mean $|\Delta|$ and Max $|\Delta|$ report the average and maximum absolute rank
displacement across the 27 countries.}
\label{tab:loo-summary}
\small
\setlength{\tabcolsep}{4pt}
\begin{tabular*}{\textwidth}{@{\extracolsep{\fill}}
  l
  S[table-format=1.4]
  S[table-format=1.4]
  S[table-format=1.2,round-precision=2]
  r
  @{}}
\toprule
{\textbf{Excluded Axis}} & {$\boldsymbol{\rho}$ \textbf{(Spear.)}} & {$\boldsymbol{\tau}$ \textbf{(Kend.)}} & {\textbf{Mean} $|\Delta|$} & {\textbf{Max} $|\Delta|$} \\
\midrule
  Financial       & 0.9927 & 0.9430 & 0.67 & 2 \\
  Energy          & 0.9921 & 0.9430 & 0.67 & 2 \\
  Technology      & 0.9683 & 0.8860 & 1.26 & 7 \\
  Defence         & 0.8327 & 0.6752 & 3.33 & 10 \\
  Critical Inputs & 0.9414 & 0.8462 & 1.78 & 8 \\
  Logistics       & 0.9261 & 0.7892 & 2.07 & 7 \\
\bottomrule
\end{tabular*}
\end{table}


\Cref{tab:loo-summary} reports the results.  Key findings:

\begin{itemize}[leftmargin=2em]
  \item \textbf{Financial axis removal} ($\rho = 0.993$).
    Dropping the financial axis barely perturbs the ranking.
    This is consistent with its low variance contribution
    (1.03\%, \cref{tab:variance-decomp}) and near-zero
    correlations with other axes.

  \item \textbf{Energy axis removal} ($\rho = 0.992$).
    Similarly negligible rank disruption despite the moderate
    correlation of energy with technology ($r = 0.53$).

  \item \textbf{Technology axis removal} ($\rho = 0.968$).
    Modest rank perturbation.  Ireland, the highest-scoring
    country on the technology axis (0.587), loses relative
    position when this axis is removed.

  \item \textbf{Defence axis removal} ($\rho = 0.833$).
    The most disruptive single omission.  Defence contributes
    35.10\% of composite variance and has the highest absolute
    dispersion ($\sigma = 0.226$).  Removing it reshuffles the
    middle ranks substantially, particularly for countries like
    the Netherlands (defence score = 0.960) and Bulgaria
    (0.861) whose high composite positions depend heavily on
    their defence concentration.

  \item \textbf{Critical inputs removal} ($\rho = 0.941$).
    Moderate disruption, consistent with its 20.92\% variance
    contribution.

  \item \textbf{Logistics axis removal} ($\rho = 0.926$).
    Moderate disruption, consistent with its 34.26\% variance
    contribution.  Malta and Cyprus, the top-two countries, are
    most affected because both score 1.00 and 0.94 respectively
    on logistics.
\end{itemize}

The full country-by-country LOO ranking table is provided in
\cref{tab:loo-detail} (\cref{app:robustness}).

\subsection{Z-score standardisation}\label{sec:zscore}

The baseline ISI uses raw (min--max bounded) axis scores.  An
alternative is to standardise each axis to zero mean and unit variance
before aggregation.  Z-score standardisation implicitly upweights
axes with higher absolute dispersion (because a one-unit change in
a high-$\sigma$ axis maps to a smaller z-increment than a one-unit
change in a low-$\sigma$ axis after rescaling).

We compute:
\begin{equation}\label{eq:zscore}
  z_{ij} = \frac{A_{ij} - \bar{A}_j}{\sigma_j}\,,
  \qquad
  \text{ISI}^{z}_i = \frac{1}{6}\sum_{j=1}^{6} z_{ij}\,,
\end{equation}
and re-rank the countries.

% Table T10: Z-score standardized composite rank comparison
\begin{table}[htbp]
\centering
\caption{Rank comparison: baseline ISI vs.\ z-score standardized composite.
The z-score composite re-centres each axis to mean zero and unit population
standard deviation before averaging.  Spearman $\rho = 0.8956$,
Kendall $\tau = 0.7265$.}
\label{tab:zscore-ranks}
\small
\setlength{\tabcolsep}{4pt}
\begin{tabular*}{\textwidth}{@{\extracolsep{\fill}}
  r l
  S[table-format=1.4]
  S[table-format=2.0,round-precision=0]
  S[table-format=+2.4]
  S[table-format=2.0,round-precision=0]
  S[table-format=2.0,round-precision=0]
  @{}}
\toprule
{\textbf{Rk}} & {\textbf{Country}} & {\textbf{Baseline}} & {\textbf{Rk\,(b)}} & {\textbf{Z-Comp.}} & {\textbf{Rk\,(z)}} & {$\boldsymbol{\Delta}$} \\
\midrule
  1 & Malta & 0.5174 & 1 & +0.8658 & 2 & 1 \\
  2 & Cyprus & 0.4682 & 2 & +0.5993 & 5 & 3 \\
  3 & Ireland & 0.4280 & 3 & +1.0732 & 1 & 2 \\
  4 & Denmark & 0.4244 & 4 & +0.6265 & 4 & 0 \\
  5 & Croatia & 0.4043 & 5 & +0.8022 & 3 & 2 \\
  6 & Finland & 0.4020 & 6 & +0.3519 & 8 & 2 \\
  7 & Sweden & 0.3987 & 7 & -0.0744 & 14 & 7 \\
  8 & Austria & 0.3824 & 8 & +0.5154 & 6 & 2 \\
  9 & Italy & 0.3738 & 9 & +0.0160 & 12 & 3 \\
  10 & Slovenia & 0.3697 & 10 & +0.2361 & 9 & 1 \\
  11 & Greece & 0.3653 & 11 & +0.1719 & 11 & 0 \\
  12 & Luxembourg & 0.3601 & 12 & -0.0826 & 15 & 3 \\
  13 & Estonia & 0.3433 & 13 & +0.4183 & 7 & 6 \\
  14 & Bulgaria & 0.3402 & 14 & -0.1427 & 18 & 4 \\
  15 & Netherlands & 0.3351 & 15 & -0.4547 & 22 & 7 \\
  16 & Hungary & 0.3296 & 16 & -0.1092 & 17 & 1 \\
  17 & Portugal & 0.3217 & 17 & -0.0251 & 13 & 4 \\
  18 & Romania & 0.3186 & 18 & +0.1849 & 10 & 8 \\
  19 & Czechia & 0.3179 & 19 & -0.1039 & 16 & 3 \\
  20 & Poland & 0.2956 & 20 & -0.4159 & 21 & 1 \\
  21 & Belgium & 0.2879 & 21 & -0.3764 & 19 & 2 \\
  22 & Germany & 0.2678 & 22 & -0.7820 & 26 & 4 \\
  23 & Lithuania & 0.2654 & 23 & -0.5512 & 23 & 0 \\
  24 & Spain & 0.2609 & 24 & -0.6720 & 24 & 0 \\
  25 & Latvia & 0.2469 & 25 & -0.6904 & 25 & 0 \\
  26 & Slovakia & 0.2364 & 26 & -0.3924 & 20 & 6 \\
  27 & France & 0.2356 & 27 & -0.9886 & 27 & 0 \\
\bottomrule
\end{tabular*}
\end{table}


\Cref{tab:zscore-ranks} compares the z-score-based ranking with the
baseline.  The Spearman rank correlation is $\rho = 0.896$---lower
than any single LOO perturbation except defence removal.  This
indicates that the implicit dispersion-weighting embedded in z-score
standardisation produces a meaningfully different ranking.

The largest rank changes under z-score standardisation include:
\begin{itemize}[leftmargin=2em]
  \item Countries with extreme scores on high-CV axes (technology,
    critical inputs) gain or lose relative position depending on
    whether the z-score upweighting favours their profile.
  \item Countries whose composite depends predominantly on the
    defence axis (which has moderate CV but the largest absolute
    $\sigma$) are less affected, because z-score standardisation
    compresses defence scores toward the mean.
\end{itemize}

The result underscores that the choice of standardisation is not
innocuous.  The ISI~v1.0 deliberately avoids z-score normalisation
because the underlying data (HHI-normalised scores) already inhabit
a common $[0, 1]$ scale, making further rescaling conceptually
unnecessary and potentially distortive.

\subsection{Dirichlet weight perturbation (Monte Carlo)}\label{sec:dirichlet}

The baseline ISI assigns equal weight ($w_j = 1/6$) to each axis.
To assess sensitivity to this choice, we draw $K = 10{,}000$
weight vectors from a symmetric $\text{Dirichlet}(\alpha = 1)$
distribution (uniform over the 5-simplex), compute the weighted
composite for each draw, rank the countries, and record the rank
distribution.

% Table T11: Dirichlet weight perturbation rank volatility
\begin{table}[htbp]
\centering
\caption{Rank volatility under Dirichlet weight perturbation
  ($\alpha_j=10$, $N_{\text{MC}}=10\,000$).}
\label{tab:dirichlet-volatility}
\small
\setlength{\tabcolsep}{4pt}
\begin{tabular*}{\textwidth}{@{\extracolsep{\fill}}
  r l
  S[table-format=2.1,round-precision=1]
  S[table-format=1.2,round-precision=2]
  S[table-format=2.0,round-precision=0]
  S[table-format=2.0,round-precision=0]
  S[table-format=3.1,round-precision=1]
  S[table-format=3.1,round-precision=1]
  @{}}
\toprule
{\textbf{Rk}} & {\textbf{Country}} & {\textbf{Mean}} & {$\boldsymbol{\sigma}$}
  & {\textbf{P5}} & {\textbf{P95}} & {\textbf{\%Top5}} & {\textbf{\%Top10}} \\
\midrule
  1 & Malta & 1.0 & 0.10 & 1 & 1 & 100.0 & 100.0 \\
  2 & Cyprus & 2.1 & 0.33 & 2 & 3 & 100.0 & 100.0 \\
  3 & Ireland & 4.0 & 1.57 & 2 & 7 & 82.0 & 99.8 \\
  4 & Denmark & 3.9 & 0.76 & 3 & 5 & 98.4 & 100.0 \\
  5 & Croatia & 5.8 & 1.75 & 3 & 9 & 44.1 & 98.2 \\
  6 & Finland & 6.4 & 2.39 & 4 & 12 & 40.6 & 93.3 \\
  7 & Sweden & 6.3 & 1.74 & 4 & 9 & 31.6 & 99.0 \\
  8 & Austria & 8.0 & 1.09 & 6 & 10 & 0.4 & 97.8 \\
  9 & Italy & 9.5 & 2.17 & 6 & 13 & 2.5 & 69.3 \\
  10 & Slovenia & 10.0 & 1.60 & 7 & 12 & 0.5 & 64.1 \\
  11 & Greece & 10.7 & 1.95 & 7 & 14 & 0.0 & 47.5 \\
  12 & Luxembourg & 11.7 & 0.99 & 10 & 13 & 0.0 & 9.4 \\
  13 & Estonia & 13.5 & 2.21 & 9 & 17 & 0.0 & 10.9 \\
  14 & Bulgaria & 14.2 & 2.03 & 11 & 18 & 0.0 & 3.8 \\
  15 & Netherlands & 15.4 & 3.07 & 10 & 20 & 0.0 & 6.7 \\
  16 & Hungary & 15.5 & 0.93 & 14 & 17 & 0.0 & 0.0 \\
  17 & Portugal & 17.0 & 1.55 & 14 & 19 & 0.0 & 0.0 \\
  18 & Romania & 17.4 & 1.65 & 15 & 19 & 0.0 & 0.1 \\
  19 & Czechia & 17.8 & 1.17 & 16 & 19 & 0.0 & 0.0 \\
  20 & Poland & 20.1 & 0.64 & 19 & 21 & 0.0 & 0.0 \\
  21 & Belgium & 20.9 & 0.56 & 20 & 22 & 0.0 & 0.0 \\
  22 & Germany & 22.9 & 1.41 & 21 & 26 & 0.0 & 0.0 \\
  23 & Lithuania & 22.9 & 0.64 & 22 & 24 & 0.0 & 0.0 \\
  24 & Spain & 23.6 & 0.75 & 22 & 24 & 0.0 & 0.0 \\
  25 & Latvia & 25.2 & 0.45 & 25 & 26 & 0.0 & 0.0 \\
  26 & Slovakia & 25.7 & 1.58 & 22 & 27 & 0.0 & 0.0 \\
  27 & France & 26.6 & 0.50 & 26 & 27 & 0.0 & 0.0 \\
\bottomrule
\end{tabular*}
\end{table}


\Cref{tab:dirichlet-volatility} reports the mean and standard
deviation of each country's rank across the 10\,000 draws, along
with the 5th and 95th percentile ranks.  Aggregate summary
statistics:

\begin{itemize}[leftmargin=2em]
  \item \textbf{Mean Spearman~$\rho$} between each draw's ranking
    and the baseline: \num{0.979} (std~=~\num{0.017}).
  \item \textbf{Minimum~$\rho$}: \num{0.863}; \textbf{Maximum~$\rho$}:
    \num{1.000}.
  \item In more than 95\% of draws, $\rho > 0.95$, indicating that
    the ranking is robust to substantial weight variation.
\end{itemize}

\textbf{Most rank-volatile countries:}
\begin{enumerate}[leftmargin=2em]
  \item Netherlands ($\sigma_{\text{rank}} = 3.07$) --- rank 15 on
    average, but ranges from 3 to 25 across draws.  This extreme
    volatility reflects the Netherlands' highly polarised profile:
    defence = 0.960 but financial = 0.121 and technology = 0.125.
  \item Finland ($\sigma_{\text{rank}} = 2.39$) --- sensitive to
    the weight placed on critical inputs (highest score: 0.547) vs
    defence (low score: 0.390).
  \item Estonia ($\sigma_{\text{rank}} = 2.21$).
  \item Italy ($\sigma_{\text{rank}} = 2.17$).
  \item Bulgaria ($\sigma_{\text{rank}} = 2.03$).
\end{enumerate}

\textbf{Most rank-stable countries:}  Malta ($\sigma_{\text{rank}} <
0.5$) and France ($\sigma_{\text{rank}} < 0.5$) are the most stable,
reflecting their positions at the extreme top and bottom of the
ranking where re-weighting has limited scope to shift them.

\subsection{Winsorization at the 95th percentile}\label{sec:winsorization}

Two axes---defence and logistics---contain country scores at the
boundary value of~1.0.  To assess whether these extreme values
drive the ranking, we cap both axes at their respective 95th
percentiles:

\begin{itemize}[leftmargin=2em]
  \item Defence: P95 cap = \num{0.94786072}.
  \item Logistics: P95 cap = \num{0.89794180}.
\end{itemize}

Any country-axis score exceeding the cap is replaced by the cap value.
The composite is then recomputed and re-ranked.

% Table T12: Winsorization sensitivity
\begin{table}[htbp]
\centering
\caption{Rank comparison: baseline vs.\ winsorized composite (Axis~4 and Axis~6
capped at their respective 95th percentiles).  Spearman $\rho = 1.0000$.}
\label{tab:winsorization}
\small
\setlength{\tabcolsep}{4pt}
\begin{tabular*}{\textwidth}{@{\extracolsep{\fill}}
  r l
  S[table-format=1.4]
  S[table-format=2.0,round-precision=0]
  S[table-format=1.4]
  S[table-format=2.0,round-precision=0]
  S[table-format=2.0,round-precision=0]
  @{}}
\toprule
{\textbf{Rk}} & {\textbf{Country}} & {\textbf{Baseline}} & {\textbf{Rk\,(b)}} & {\textbf{Winsor.}} & {\textbf{Rk\,(w)}} & {$\boldsymbol{\Delta}$} \\
\midrule
  1 & Malta & 0.5174 & 1 & 0.4917 & 1 & 0 \\
  2 & Cyprus & 0.4682 & 2 & 0.4610 & 2 & 0 \\
  3 & Ireland & 0.4280 & 3 & 0.4280 & 3 & 0 \\
  4 & Denmark & 0.4244 & 4 & 0.4244 & 4 & 0 \\
  5 & Croatia & 0.4043 & 5 & 0.4043 & 5 & 0 \\
  6 & Finland & 0.4020 & 6 & 0.4020 & 6 & 0 \\
  7 & Sweden & 0.3987 & 7 & 0.3987 & 7 & 0 \\
  8 & Austria & 0.3824 & 8 & 0.3824 & 8 & 0 \\
  9 & Italy & 0.3738 & 9 & 0.3738 & 9 & 0 \\
  10 & Slovenia & 0.3697 & 10 & 0.3697 & 10 & 0 \\
  11 & Greece & 0.3653 & 11 & 0.3653 & 11 & 0 \\
  12 & Luxembourg & 0.3601 & 12 & 0.3601 & 12 & 0 \\
  13 & Estonia & 0.3433 & 13 & 0.3433 & 13 & 0 \\
  14 & Bulgaria & 0.3402 & 14 & 0.3402 & 14 & 0 \\
  15 & Netherlands & 0.3351 & 15 & 0.3332 & 15 & 0 \\
  16 & Hungary & 0.3296 & 16 & 0.3296 & 16 & 0 \\
  17 & Portugal & 0.3217 & 17 & 0.3217 & 17 & 0 \\
  18 & Romania & 0.3186 & 18 & 0.3186 & 18 & 0 \\
  19 & Czechia & 0.3179 & 19 & 0.3179 & 19 & 0 \\
  20 & Poland & 0.2956 & 20 & 0.2956 & 20 & 0 \\
  21 & Belgium & 0.2879 & 21 & 0.2879 & 21 & 0 \\
  22 & Germany & 0.2678 & 22 & 0.2678 & 22 & 0 \\
  23 & Lithuania & 0.2654 & 23 & 0.2654 & 23 & 0 \\
  24 & Spain & 0.2609 & 24 & 0.2609 & 24 & 0 \\
  25 & Latvia & 0.2469 & 25 & 0.2469 & 25 & 0 \\
  26 & Slovakia & 0.2364 & 26 & 0.2364 & 26 & 0 \\
  27 & France & 0.2356 & 27 & 0.2356 & 27 & 0 \\
\bottomrule
\end{tabular*}
\end{table}


\Cref{tab:winsorization} compares the winsorized ranking with the
baseline.  The Spearman rank correlation is $\rho = 1.000$: the
ranking is entirely preserved.  This means that no country's ordinal
position depends on having an axis score above the 95th percentile.
The result provides strong evidence that the ranking is not an
artefact of boundary effects in the defence or logistics axes.

\subsection{Robustness summary}

\Cref{tab:robustness-summary} consolidates the four robustness
exercises.

\begin{table}[htbp]
\centering
\caption{Robustness programme summary.  Spearman~$\rho$ between
  each perturbed ranking and the ISI~v1.0 baseline.}
\label{tab:robustness-summary}
\begin{tabular}{@{}llS[table-format=1.3]l@{}}
\toprule
\textbf{Exercise} & \textbf{Perturbation} & {$\boldsymbol{\rho}$}
  & \textbf{Verdict} \\
\midrule
LOO (Financial)   & Drop axis 1 & 0.993 & Stable \\
LOO (Energy)      & Drop axis 2 & 0.992 & Stable \\
LOO (Technology)  & Drop axis 3 & 0.968 & Stable \\
LOO (Defence)     & Drop axis 4 & 0.833 & Moderate sensitivity \\
LOO (Critical)    & Drop axis 5 & 0.941 & Stable \\
LOO (Logistics)   & Drop axis 6 & 0.926 & Stable \\
Z-score           & Re-standardise & 0.896 & Moderate sensitivity \\
Dirichlet MC      & Random weights & 0.979 & Stable (mean) \\
Winsorization     & Cap at P95 & 1.000 & Fully stable \\
\bottomrule
\end{tabular}
\end{table}

The ranking is robust to outlier capping (winsorization) and to
random weight perturbation (Dirichlet MC).  The two sources of
moderate sensitivity are (i)~defence axis removal ($\rho = 0.833$)
and (ii)~z-score standardisation ($\rho = 0.896$).  Both of these
are structural: defence is the single most dispersed and
variance-contributing axis, and z-score standardisation introduces
implicit dispersion-based reweighting.  Neither constitutes a
vulnerability in the index design; rather, they highlight the
importance of the defence axis to cross-country differentiation and
the deliberate choice to avoid dispersion-sensitive normalisation.
